\section{结论}\label{sec:conclusion}

本文综述了黎曼猜想从1859年提出至今的研究历程。可以概括出三条主线：

\begin{enumerate}[label=(\arabic*)]
    \item \textbf{解析主线}：从黎曼的显式公式到哈代—李特尔伍德—塞尔伯格的硬性定理，再到当前 $41.7\%$ 的临界线零点比例，解析数论在“平均意义”上不断逼近真理，但始终无法排除个别离群零点。
    \item \textbf{几何主线}：韦伊与德利涅在有限域上的成功证明了黎曼型猜想在正确框架下是可以解决的，这为朗兰兹纲领与未来可能的上同调理论指明了方向。
    \item \textbf{物理与计算主线}：随机矩阵理论、希尔伯特—韦尔纲领与大规模数值验证从经验层面提供了压倒性的支持，并将猜想与量子混沌、谱理论深刻联系起来。
\end{enumerate}

黎曼猜想的意义早已超出素数分布本身：它是检验数学统一性的试金石，其任何解决——无论是证明还是反例——都将引发数论、代数几何与数学物理的连锁变革。正如蒙哥马利与沃尔克（Vaughan）所言，我们拥有“证据的海洋”，但缺少的恰恰是一座桥。未来的突破更可能来自新结构的发现（正确的上同调理论、迹公式或谱实现），而非现有方法的渐进改良。
